\begin{abstract}
Light carries data and signatures of the path and objects it encounters,
allowing the same optical infrastructure to support communication and sensing.
Optical integrated sensing and communication (O-ISAC) builds on this dual
capability across fiber, free space optical, visible light and light fidelity
(LiFi), millimeter
wave and terahertz systems enabled by photonics, and hybrid systems. This
breadth raises a central question: which design insights are shared, and which
depend on a particular platform? A systematic literature search covering
1 January 2020 through 22 June 2026 identified 227 eligible reports,
corresponding to 206 unique studies. Using structured narrative synthesis and a
review specific
technical assessment, we organized the evidence by physical modality and by
how communication and sensing are coupled. The review received no specific
support and was retrospectively registered on OSF (7f6wb; DOI:
10.17605/OSF.IO/7F6WB). Our comparison framework keeps each result connected to
its task, measurement plane, operating conditions, and validation setting. The
evidence reveals conditional relations in which bandwidth, optical power,
hardware constraints, and signal design jointly shape communication reliability
and sensing performance, extending design reasoning beyond a single rate and
resolution frontier. Laboratory demonstrations already span a wide range of
architectures; field validation, shared benchmarks, and reusable data and code
offer the next stage of maturity. Together, the taxonomy and framework help
researchers interpret evidence, choose meaningful benchmarks, and advance
O-ISAC toward 6G.
\end{abstract}

\begin{IEEEkeywords}
6G, integrated sensing and communication, metric comparability, optical
communication, optical sensing.
\end{IEEEkeywords}
